% ===== AAAI-27 spine draft: Introduction (~0.9 col-page target) =====
\section{Introduction}\label{sec:intro}

Between reading ``turn the car red'' and painting it, what does a
diffusion image editor actually do? The mechanics of diffusion suggest
a gradual-development story, each of sixty blocks and twenty denoising
steps contributing its small share -- but this paper measures that
story and finds it wrong. In the model we study, the realized edit is
\emph{predictable early} from an internal target-image code, and only
\emph{translated late}, in the last few blocks, into the flow field that does
the painting.

Prior work provides partial readouts of generative internals. The logit and
tuned lenses project transformer residuals through a final or learned output
map~\citep{nostalgebraist2020logitlens,belrose2023tuned}. In diffusion,
Diffusion Lens interprets text-encoder states; DiffLens and ICM probe or steer
denoiser features; and DreamReader and sparse autoencoders provide broader
extraction and feature-analysis tools
~\citep{toker2024diffusionlens,shi2025difflens,zaleska2026icm,prakash2026dreamreader,surkov2025sae}.
Linear-representation work motivates a probe-readable basis, while multimodal
reasoners externalize visual thoughts as sketches or whiteboards
~\citep{park2024linear,hu2024sketchpad,menon2024whiteboard,wu2024vot,li2025mvot};
we instead test for an unrendered internal image, complementing the broader
thinking-with-images taxonomy~\citep{su2025thinkingsurvey}.
Recent work also detects early cross-seed convergence, predicts final quality
from early activations, or allocates adaptive image-editing budgets
~\citep{kwon2026lockin,guo2026probeselect,qu2026adecot}. None jointly decodes
an instruction-based DiT's image stream across both transformer depth and
denoising time while separating explicit-prompt readability from the realized
outcome under a fixed prompt.

A second line studies how to produce edits or intervene on hidden states.
Modern editors build on diffusion, latent/flow formulations, and guidance
~\citep{ho2020ddpm,song2021ddim,rombach2022ldm,lipman2023flowmatching,liu2023rectifiedflow,ho2022cfg}.
InstructPix2Pix and Prompt-to-Prompt learn or manipulate diffusion editing
controls~\citep{brooks2023instructpix2pix,hertz2023prompttoprompt}; ROME and
activation-patching practice transfer hidden states, while activation
engineering adds steering directions
~\citep{meng2022rome,heimersheim2024activation,turner2023actadd}. Concept
Lancet already uses representation transplant for diffusion editing
~\citep{luo2025colan}. Recent analyses further show that detectable diffusion
features can fail under steering and that multi-mediator patching can hide
interaction effects~\citep{cassano2026look,vaidyanathan2026curse}. We therefore
use transplant only as a sufficiency test and steering only as a bounded
writability test, not as a standalone editing method or a necessity proof.

We build a layer--time lens for a production DiT editor,
\model{} (60 blocks)~\citep{peebles2023dit,qwen2025image}, by decoding intermediate residual streams through
the model's own frozen output head, and -- crucially -- pair passive
readings with linear probes and controlled activation
interventions (ablation, transplant, steering) that test what those
activations can install or write.

Naively porting the logit-lens recipe reproduces the familiar mid-stack
``nothing happens yet'' appearance -- and that appearance is itself
misleading. With source and an underspecified recoloring instruction
fixed across 48 seeds, a generation-level probe predicts the color
ultimately chosen from layer 6, where the frozen head still reads noise.
The outcome is held in a basis the head cannot parse until a narrow band
near the top translates it into the velocity field that paints the result.

The main contributions of this work are as follows:
\begin{itemize}
\item \textbf{An instrument.} A layer--time instrument for a DiT image
editor, combining a training-free frozen-head lens with tuned and
differential variants,
paired with linear probes and three intervention primitives (ablate,
transplant, steer) with stated write semantics, under a
falsification-first protocol -- pre-registered predictions, per-image
visual verdicts, and bit-exact gates.
\item \textbf{A mechanism: the editor thinks with an internal image.}
The realized edit is predictable early under a fixed underspecified
instruction, carried as a target-image code, and translated into the
velocity code near layers 52--54. Across the two densely measured
timesteps, the crossover midpoint differs by under one layer while its
width varies from one to four layers; the boundary also replicates in a
same-lineage fine-tune and a 4-step distillation. Transplant establishes
sufficiency; in the tested recolor setting, steering and a rank ladder show
that readable attributes require a spatial carrier and a
variable-dependent write subspace.
\item \textbf{Validation and a lever.} Held-out early best-of-five selection
validates the readout on Qwen and an independent architecture, with about
68\% measured compute reduction. Separately, truncating classifier-free
guidance after step 2 saves 45\% of forward passes across tested edit
families; dropping CFG entirely produces failures consistent with a
non-redundant role in scope resolution rather than outcome selection.
\end{itemize}
This is exploratory, single-model-family work; every verdict was made
on the images themselves, and we are explicit about where the lens is
unreliable.
